\documentclass[11pt,a4paper]{article}
\usepackage[T1]{fontenc}
\usepackage[utf8]{inputenc}
\usepackage{newtxtext}
\usepackage[margin=1in]{geometry}
\usepackage{enumitem}
\usepackage[hidelinks]{hyperref}

\hypersetup{
  pdftitle={Prior Publication Statement -- MEET-COBRA},
  pdfauthor={Bowen Xie, Sheng Zhou, and Zhisheng Niu},
  pdfsubject={Relationship to the AEPHORA conference paper presented at IEEE ICC 2025}
}
\pagestyle{empty}
\setlength{\parindent}{0pt}
\setlength{\parskip}{0.55em}

\begin{document}

\begin{center}
  {\Large\bfseries Prior Publication Statement}\par
  \vspace{0.4em}
  {\bfseries MEET-COBRA: Machine-learning-based Energy-Efficient proacTive
  Coordination of handOver, Beamforming and Resource Allocation}\par
  Bowen Xie, Sheng Zhou, and Zhisheng Niu\par
  Manuscript ID: Paper-TW-Feb-26-0360\par
  Submitted to \emph{IEEE Transactions on Wireless Communications}
\end{center}

\textbf{Earlier publication}

Bowen Xie, Sheng Zhou, Zhisheng Niu, Hao Wu, and Cong Shi,
``AEPHORA: AI/ML-Based Energy-Efficient Proactive Handover and Resource Allocation,''
in \emph{ICC 2025 -- IEEE International Conference on Communications},
Montreal, Canada, June 8--12, 2025, pp.~6297--6303.\\
DOI: \href{https://doi.org/10.1109/ICC52391.2025.11161541}{10.1109/ICC52391.2025.11161541}.
The conference was sponsored by the IEEE Communications Society
(\href{https://icc2025.ieee-icc.org/}{conference website}).

\textbf{Relationship and main extensions}

The journal manuscript extends the above conference paper. Both works consider
a two-tier vehicular heterogeneous network and use prediction-assisted proactive
handover (PHO) and queue-aware resource allocation (RA) to reduce system transmit
power under latency constraints. The main extensions are as follows.

\begin{enumerate}[leftmargin=*,itemsep=0.55em,parsep=0pt,topsep=0.2em]
  \item \textbf{Joint coordination of PHO, beamforming, and RA.}
  The conference paper uses a scalar channel model with a prescribed beamforming
  gain and does not optimize beam selection. The journal manuscript introduces
  microcell multiple-input multiple-output (MIMO) channels, transmit and receive
  beam-pair selection, and beam-training overhead into the joint optimization.
  It develops prediction-guided beam search with early termination and coordinates
  it with frame-level PHO and slot-level RA. It also accounts for inter-cell
  interference and handover interruption when estimating resource demands and allocating
  resources.

  \item \textbf{Prediction of beam pairs and channel gains.}
  The conference paper predicts vehicle positions from historical channel gains
  and estimates future link gains from the predicted positions. The journal
  manuscript instead uses historical superposed-pilot observations to predict
  next-frame beam-pair candidates, desired-link gains, and beam-averaged
  interfering-link gains. These predictions guide association decisions, beam
  search, and resource allocation without requiring full MIMO channel reporting.

  \item \textbf{Ray-traced channels and broader evaluation.}
  The conference evaluation combines SUMO vehicle trajectories with statistically
  generated scalar channels. The journal evaluation uses SUMO mobility traces and
  Sionna Ray Tracing to generate microcell MIMO channels, and computes receive
  rates using the selected beams and explicit resource block assignments. It
  includes comparisons with two adapted joint handover and beamforming methods from
  the literature, module ablations, multiple random seeds, analyses of mobility
  and prediction errors, and measurements of neural-network complexity,
  inference latency, and reporting overhead.
\end{enumerate}

The conference paper is cited as reference~[13] in the journal manuscript.
Section~I, page~2, explicitly identifies the manuscript as an extension and
summarizes the principal differences. A copy of the conference paper is provided
separately with this statement.

\end{document}
